Orbiter has an interface to GAP~\cite{GAP}. 
This allows Orbiter commands to utilize GAP specific functionality. 
One application is the computation of canonical images 
of sets in the context of permutation groups. 
This is facilitated by means of the GAP package images~\cite{images}.

\bigskip


Table~\ref{tab:grouptasks:GAP} 
lists the group theoretic commands in Orbiter 
that rely on GAP.
\orbitertableofcommands{group_theoretic_activity_gap}

\bigskip


How can we export a group from Orbiter to GAP? 
For semilinear matrix groups, 
we need to export to the GAP package Fining~\cite{fining}. 
This is because semilinear groups are not available in GAP.
Here is an example. 
We create the stabilizer of the Hirschfeld surface 
in $\PG(3,4)$ and export the group to GAP/Fining 
as a group of collineations.

\sourcecode{Hirschfeld_stab_export_gap}

The command creates a GAP/Fining file shown below:
\orbiteroutput{GRAPHICS/LATEX/Hirschfeld_stab_fining}

When we load this file into GAP, we see the following output
(for printing purposes, 
it has been shortened and reformatted slightly):
\orbiteroutput{GRAPHICS/LATEX/Hirschfeld_stab_fining_gap_output}

\bigskip



Suppose we want to compute the least image 
of the set $\{3,5\}$ in the permutation group 
$C_6.$
The following command sets up the group 
$C_6$ and invokes the Orbiter \verb'-canonical_image' command.

\sourcecode{Cyclic_6_canonical_image}

The command creates a GAP file shown below:
\orbiteroutput{GRAPHICS/LATEX/minimal_image_cyclic_6_command}
Running this command in GAP yields:
\orbiteroutput{GRAPHICS/LATEX/minimal_image_cyclic_6_gap_output}
Note that the permutation representations in GAP are 1-based, 
while Orbiter uses 0-based representations. 
So, the least image of $\{1,3\}$ 
in GAP really corresponds to the set $\{0,2\}$ in Orbiter.


\bigskip


Let us try a larger example. 
We want to compute the canonical form of the 
Edge curve by means of their rational points over $\bbF_{17}.$
We use the following makefile variable 
to specify the set of rational points:

\sourcecodevariable{EDGE_CURVE_Q17_AS_POINTS}

The next command creates $\PGL(3,17)$ 
and invokes the \verb'-canonical_image' image command.

\sourcecode{Edge_qurve_q17_canonical_image}

The command creates a GAP input file called \verb'PGL_3_17_canonical_image.gap'. 
The file is too large to be shown here.
%\orbiteroutput{GRAPHICS/LATEX/minimal_image_PSL_3_17_command}
Running it through GAP yields the canonical form:
\orbiteroutput{GRAPHICS/LATEX/minimal_image_PSL_3_17_gap_output}
Keep in mind that the permutation representation in Gap is 1-based, 
so the numbers have been incremented by one.
